\scanpage{55}
\[
L_{E_\eta}\left(\frac1{qt}\right)=A(t)L_{E_\eta}(t)
\]

\[
A(t)=\prod_x\tikz[remember picture,baseline=(localfactor.base)]{\node[draw,ellipse,inner sep=2pt] (localfactor) {$\displaystyle\frac{L_{\alpha_x(E_\eta)}(t^{-1})}{L_{\alpha_x(E_\eta)}(p^{-\rho}t)}$};}
(-t)^{\chi^{\unsure{\leftrightharpoons|\leftrightarrows|\natural}}(E_\eta)-\sum_x\chi_x(E_\eta)}
\frac{q^{\chi^{\leftrightharpoons}(E_\eta)}}{\delta^{\leftrightharpoons}(E_\eta)\prod_x\delta_x(E_\eta)^{-1}}
\]
\hspace{.42\linewidth}\tikz[remember picture,baseline=(factorlabel.base)]{\node[inner sep=0pt] (factorlabel) {$\inserted{=A_x(t^{-1})^{-1}}$};}
\tikz[remember picture,overlay]{\draw (localfactor.south east) to[out=-20,in=180] (factorlabel.west);}

\[
L_{\alpha_x(\deleted{\unsure{E}}E_\eta)}\left(\frac1{p^\rho t}\right)=A_x(t)L_{\alpha_x(E_\eta)}(t)
\]
\[
\deleted{L_{\alpha_x(\unsure{E_\eta})}\quad L_{\alpha_x(E_\eta)}(}
\]

N.B. Ayant coupé $\lambda_x(\deleted{\unsure{E_\eta}}E_\eta)(t)$ \deleted{\unsure{en}} en morceaux, on en perd le \uline{\unsure{caractère} additif}, Il vaut sans doute mieux ne pas faire ce découpage ---

\uline{Pb} a) Dans les cas géométriques \unsure{classiques}, construire les $\alpha_x(E_\eta)$ qui \unsure{déterminent} les $\mu_x$, donc les $\lambda_x(E_\eta)(t)$. Déterminer en particulier \tikz[remember picture,baseline=(chiloc.base)]{\node[inner sep=0pt] (chiloc) {$\chi_x(E_\eta)$};}\margin{\tikz[remember picture,baseline=(geolabel.base)]{\node[inner sep=0pt] (geolabel) {\shortstack{\unsure{invariant}\\\uuline{géom.}}};}}, \tikz[remember picture,baseline=(deltaloc.base)]{\node[inner sep=0pt] (deltaloc) {$\delta_x(E_\eta)$};}\margin{\tikz[remember picture,baseline=(arithlabel.base)]{\node[inner sep=0pt] (arithlabel) {\shortstack{\unsure{invariant}\\\uuline{arithm.}}};}}, les \unsure{valeurs absolues} des val propres de \unsure{frobenius|Frobenius}.

b) Donner une formule genre \unsure{Chafarevitch|Chafarévitch} - Ogg pour \tikz[remember picture,baseline=(chiglob.base)]{\node[inner sep=0pt] (chiglob) {$\chi^{\leftrightharpoons}(E_\eta)$};}

c) Donner une formule pour \tikz[remember picture,baseline=(deltaglob.base)]{\node[inner sep=0pt] (deltaglob) {$\delta^{\leftrightharpoons}(E_\eta)$};}
\tikz[remember picture,overlay]{
\draw[->] (chiloc.south) to[out=-85,in=155] (geolabel.west);
\draw[->] (chiglob.east) to[out=0,in=-105] (geolabel.south);
\draw[->] (deltaloc.south) to[out=-80,in=100] (arithlabel.north);
\draw[->] (deltaglob.east) to[out=0,in=-135] (arithlabel.south);
}
% archive sheet number: 54
